\title{DeepSeekMath: Pushing the Limits of Mathematical Reasoning in Open Language Models}

\begin{abstract}
The task of mathematical reasoning remains particularly difficult for language models, owing to its inherently intricate and highly organized structure. This study presents DeepSeekMath~7B, a model developed by continuing the pre-training of DeepSeek-Coder-Base-v1.5 7B with an extensive dataset comprising 120B math-centric tokens derived from Common Crawl, alongside datasets of natural language and code. DeepSeekMath~7B attains a notable 51.7\% score on the rigorous, competition-level MATH benchmark, without leveraging additional toolkits or ensemble (voting) strategies, and its performance approaches that of Gemini-Ultra and GPT-4. When applying self-consistency with 64 generated samples, DeepSeekMath~7B reaches a score of 60.9\% on the MATH benchmark. The robust mathematical reasoning skills exhibited by DeepSeekMath~stem from two central innovations: first, a carefully crafted data curation pipeline that makes full use of public web resources; and second, the introduction of Group Relative Policy Optimization (GRPO)—a refinement of Proximal Policy Optimization (PPO)—which boosts both reasoning proficiency and the efficiency of PPO in memory usage.
\end{abstract}

\section{Introduction}

Large language models (LLM) have transformed methods of mathematical reasoning within artificial intelligence, leading to notable progress on benchmarks for quantitative reasoning \citep{MATH} and geometry reasoning \citep{trinh2024solving}. In addition, these systems have shown themselves to be valuable tools for aiding humans in tackling intricate mathematical challenges \citep{tao}. Nevertheless, advanced models like GPT-4 \citep{gpt4} and Gemini-Ultra \citep{gemini} remain inaccessible to the public, while open-source alternatives currently available lag substantially in terms of their capabilities.

In this work, we present DeepSeekMath, a specialized language model tailored for mathematics that markedly surpasses open-source alternatives and approaches the performance of GPT-4 on standard academic benchmarks. This advancement is underpinned by the construction of the DeepSeekMath Corpus—a massive, high-quality pre-training dataset that consists of 120B mathematical tokens. The data for this corpus is collected from Common Crawl (CC) with the aid of a classifier based on fastText \citep{joulin2016fasttext}. During the initial stage, the classifier is trained using examples from OpenWebMath \citep{openwebmath} as positive samples and a broad array of unrelated web pages as negative samples. Once trained, the classifier is used to extract additional positive samples from CC, which are subsequently curated with human annotation. The classifier is then retrained with these improved datasets to further enhance its accuracy. Evaluation demonstrates that the resulting large-scale corpus is of superior quality, with our foundational model, DeepSeekMath-Base 7B, reaching 64.2\% accuracy on GSM8K \citep{gsm8k} and 36.2\% on the advanced MATH benchmark \citep{MATH}, surpassing the performance of Minerva 540B \citep{minerva}. Moreover, as the DeepSeekMath Corpus incorporates multiple languages, we observe notable performance gains on Chinese mathematics benchmarks \citep{wei2023cmath,agieval}. We regard our methodology for processing mathematical data as a useful foundation for further exploration by the research community and acknowledge that considerable advancements remain possible.

DeepSeekMath-Base is built upon DeepSeek-Coder-Base-v1.5 7B \citep{deepseek-coder}, as our findings suggest that initializing from a code-specialized model yields superior results over beginning with a standard large language model. Additionally, our experiments reveal that math-focused training not only elevates performance on MMLU \citep{mmlu} and BBH benchmarks \citep{bbh}, but also broadly strengthens the model’s reasoning skills beyond mathematics.

Following pre-training, we perform mathematical instruction tuning on DeepSeekMath-Base utilizing data featuring chain-of-thought \citep{cot}, program-of-thought \citep{pot,pal}, and tool-integrated reasoning \citep{tora} methods. The tuned model, DeepSeekMath-Instruct 7B, surpasses all other 7B models and demonstrates performance on par with open-source instruction-tuned models at the 70B scale.

In addition, we propose Group Relative Policy Optimization (GRPO), which is a modified reinforcement learning (RL) algorithm based on Proximal Policy Optimization (PPO) \citep{schulman2017proximal}. Unlike PPO, GRPO eliminates the need for a critic network and instead derives its baseline from aggregated group scores, thereby considerably lowering computational requirements during training. Using only a portion of English instruction tuning data, GRPO yields significant gains compared to the robust DeepSeekMath-Instruct across both in-domain benchmarks (GSM8K: 82.9\% $\rightarrow$ 88.2\%, MATH: 46.8\% $\rightarrow$ 51.7\%) as well as on out-of-domain problems (e.g., CMATH: 84.6\% $\rightarrow$ 88.8\%) in the reinforcement learning phase. We further develop a unified framework that relates various approaches, including Rejection Sampling Fine-Tuning (RFT) \citep{yuan2023scaling}, Direct Preference Optimization (DPO) \citep{dpo}, PPO, and GRPO. Within this framework, it becomes clear that these approaches can all be interpreted as forms of direct or simplified reinforcement learning. Comprehensive experiments are performed to thoroughly analyze key factors of this framework, exploring, for example, online versus offline learning, supervision over outputs versus processes, and comparisons between single-step and iterative RL. Finally, we discuss how our RL methodology enhances the performance of instruction-tuned models and outline prospective research avenues for advancing RL strategies within this unified conceptual framework.

\subsection{Contributions}
We present a scalable approach to math pre-training and provide a comprehensive investigation and evaluation of reinforcement learning methods.

\textbf{Math Pre-Training at Scale}

\begin{itemize}[topsep=0pt]
    \item This study demonstrates that the Common Crawl dataset, which is openly available, holds significant mathematical information. Through the use of a carefully crafted data selection framework, we assemble the DeepSeekMath Corpus—a high-quality dataset comprising 120B tokens extracted from mathematically relevant web pages. This corpus is nearly sevenfold larger than the collection of math-related web content employed by Minerva \citep{minerva}, and surpasses the recently introduced OpenWebMath \citep{openwebmath} by a factor of nine.

\item DeepSeekMath-Base 7B, our foundational pre-trained model, matches the performance of Minerva 540B \citep{minerva}, demonstrating that model size alone does not determine proficiency in mathematical reasoning. High-quality pre-training data can enable smaller models to exhibit robust performance.

\item We present results from our experiments focused on math training. Introducing code training before engaging in math tasks enhances the model's performance on mathematical problems, regardless of whether external tools are used. These outcomes contribute to the ongoing discussion around the question: \textit{does code training enhance reasoning skills?} Our evidence suggests that it does, particularly when it comes to mathematical reasoning.

\item While utilizing arXiv papers for training is a widespread practice, particularly among mathematics-focused research, it does not yield significant gains on any of the mathematical benchmarks considered in this study.
\end{itemize}

\textbf{Exploration and Analysis of Reinforcement Learning}

\begin{itemize}[topsep=0pt]
    \item We propose Group Relative Policy Optimization (GRPO), a novel reinforcement learning method that eliminates the need for a critic model by deriving the baseline from group scores. This design choice results in substantial savings in computational resources when compared to Proximal Policy Optimization (PPO).
    \item We show that applying GRPO markedly improves the performance of our instruction-tuned model, DeepSeekMath-Instruct, utilizing only instruction-tuning data. Additionally, we find that reinforcement learning with GRPO leads to improvements in generalization to out-of-domain data.
    \item We present a unified framework to analyze and contextualize a variety of approaches, including RFT, DPO, PPO, and GRPO. Our work includes comprehensive experimental analyses—such as comparing online versus offline training, assessing outcome versus process supervision, and contrasting single-step with iterative reinforcement learning—to thoroughly examine the critical components underlying this unified framework.
    \item Building on our unified perspective, we investigate the fundamental factors contributing to the effectiveness of reinforcement learning and outline several promising research avenues for optimizing reinforcement learning in large language models (LLMs).
\end{itemize}

\subsection{Summary of Evaluations and Metrics}

\begin{itemize}[topsep=0pt]
    \item \textbf{English and Chinese Mathematical Reasoning}: Our evaluation encompasses a broad array of English and Chinese benchmarks, with tasks ranging from elementary to collegiate mathematics. The English datasets utilized comprise GSM8K \citep{gsm8k}, MATH \citep{MATH}, SAT \citep{llemma}, OCW Courses \citep{minerva}, and MMLU-STEM \citep{mmlu}. For Chinese, assessments are conducted with MGSM-zh \citep{mgsm}, CMATH \citep{wei2023cmath}, Gaokao-MathCloze \citep{agieval}, and Gaokao-MathQA \citep{agieval}. Our analysis measures the models' proficiency both in generating comprehensive, stand-alone written solutions without assistance from external tools, and in problem solving through Python code.

On English evaluation benchmarks, DeepSeekMath-Base matches the performance of the proprietary Minerva 540B model \citep{minerva} and outperforms all open-source base models—including Mistral 7B \citep{mistral} and Llemma-34B \citep{llemma}—regardless of prior math-specific pre-training, often by a wide margin. Importantly, DeepSeekMath-Base also establishes state-of-the-art results on Chinese mathematical benchmarks. This advantage likely arises from our divergence from earlier works \citep{minerva,llemma}, as we augment the training corpus with both English and high-quality non-English mathematical data. Following mathematical instruction tuning and reinforcement learning, the derived DeepSeekMath-Instruct and DeepSeekMath-RL models demonstrate robust capabilities, for the first time allowing an open-source model to surpass 50\% accuracy on the challenging competition-level MATH dataset.

\item \textbf{Formal Mathematics}: We assess DeepSeekMath-Base through the informal-to-formal theorem proving benchmark introduced in \citep{dsp_proof}, utilizing miniF2F \citep{minif2f} and selecting Isabelle \citep{isabelle} as the proof assistant for this task. DeepSeekMath-Base achieves impressive few-shot autoformalization results.
\item \textbf{Natural Language Understanding, Reasoning, and Code}: For a well-rounded evaluation of the model’s ability in language comprehension, reasoning, and programming, we test DeepSeekMath-Base using the Massive Multitask Language Understanding (MMLU) benchmark \citep{mmlu}, which includes 57 varied multiple-choice disciplines; the BIG-Bench Hard (BBH) suite \citep{bbh} comprising 23 complex, multi-step reasoning challenges; and the coding evaluation datasets HumanEval \citep{codex} and MBPP \citep{mbpp}. Pre-training on mathematical data enhances the model’s performance across understanding, reasoning, and coding competencies.

\section{Math Pre-Training}

\subsection{Data Collection and Decontamination} 

Here, we describe the methodology used to build the DeepSeekMath Corpus utilizing Common Crawl. Figure \ref{fig:data_collect} illustrates our stepwise pipeline designed for comprehensive collection of a large mathematical corpus, initiating with a seed corpus—namely, a limited yet carefully curated selection of math-related datasets. Importantly, this methodology can be extended for use in other fields, including code datasets.

We begin by selecting OpenWebMath \citep{openwebmath}, a curated dataset of mathematical web documents, as our foundational seed corpus. Leveraging this dataset, we construct a fastText model \citep{joulin2016fasttext} to identify additional web pages with similar mathematical content. For model training, we randomly extract 500,000 examples from the seed corpus to serve as positive instances and pair these with 500,000 negative examples drawn from Common Crawl. Model training is conducted using an open-source tool\footnote{\url{https://fasttext.cc}}, with hyperparameters set as follows: a vector dimension of 256, learning rate of 0.1, maximum n-gram length of 3, minimum word occurrence threshold of 3, and training over 3 epochs. To make the size of Common Crawl manageable, we apply deduplication strategies based on URLs and near-duplicate detection, yielding a dataset of 40 billion unique HTML web pages. Utilizing the trained fastText model, we retrieve candidate mathematical documents from the deduplicated Common Crawl set. To exclude content of lower quality, these retrieved pages are ranked by their predicted fastText scores, and only the highest-scoring documents are retained. We evaluate how much data to keep by pre-training on subsets with the top 40B, 80B, 120B, and 160B tokens. For the initial pass, we opt to retain only the top 40B tokens.

Following the initial round of data acquisition, a significant number of mathematical web pages were not gathered, primarily because the fastText model had been trained on a seed set of positive samples that did not encompass sufficient variation. To address this limitation, we seek out and incorporate additional mathematical web sources to expand the initial corpus, thereby refining the fastText model. Concretely, we begin by dividing the entire Common Crawl dataset into separate domains, defining each domain as the collection of web pages that share the same base URL. For every domain, we compute the proportion of pages captured during the first collection iteration. If a domain exhibits over 10\% of its pages having been gathered, we consider it math-related (for instance, \url{mathoverflow.net}).
Next, we conduct manual annotation of URLs containing mathematical material within the identified domains (e.g., \url{mathoverflow.net/questions}).
Any web pages associated with these URLs that were missed in previous rounds are incorporated into the growing seed corpus. This process enables the inclusion of a broader array of positive samples, which in turn allows the fastText model to learn from more comprehensive mathematical data and improves its coverage in subsequent rounds of data collection. Across four rounds, this methodology yields 35.5M mathematical web pages, amassing a total of 120B tokens. By the fourth collection, we observe that approximately 98\% of the new data had already been obtained during the third round, leading us to conclude the collection process.

To prevent the issue of benchmark leakage, we adopt the filtering strategy described by \cite{deepseek-coder}, excluding web content that includes questions or answers from established English mathematical datasets like GSM8K \citep{gsm8k} and MATH \citep{MATH}, as well as Chinese resources such as CMATH \citep{wei2023cmath} and AGIEval \citep{agieval}. The filtration process operates as follows: any text fragment within our math training data that precisely matches a 10-gram sequence found in the benchmark datasets is eliminated. Additionally, for benchmark entries between 3 and 9 grams, we remove web pages based on exact string correspondence to further ensure the integrity of our dataset.

\subsection{Validating the Quality of the DeepSeekMath~Corpus}

We run pre-training experiments to investigate how the DeepSeekMath~Corpus is compared with the recently released math-training corpora:

\begin{itemize}[topsep=0pt]
    \item \textbf{MathPile} \citep{mathpile}: This dataset is comprised of 8.9 billion tokens collected from a range of sources, including textbooks, Wikipedia, ProofWiki, CommonCrawl, StackExchange, and arXiv. Notably, arXiv constitutes over 85\% of the total tokens;
    \item \textbf{OpenWebMath} \citep{openwebmath}: Extracted from CommonCrawl and carefully filtered for mathematical material, this resource contains 13.6 billion tokens;
    \item \textbf{Proof-Pile-2} \citep{llemma}: This large mathematical dataset brings together OpenWebMath, AlgebraicStack (which includes 10.3 billion tokens of mathematical code), and 28.0 billion tokens from arXiv publications. In line with \cite{llemma}, we adopt an arXiv:Web:Code composition ratio of 2:4:1 in our experiments using Proof-Pile-2.
\end{itemize}

\subsubsection{Training Setting}
\label{sec:quality-policy}
Mathematical training is performed on a general-purpose language model with 1.3B parameters, utilizing the architecture of the DeepSeek LLMs \citep{deepseek-llm}, here referred to as DeepSeek-LLM 1.3B. Each mathematical dataset serves as the basis for separately training a corresponding model for a total of 150B tokens. All experimental procedures employ the HAI-LLM \citep{haillm} training infrastructure, known for its efficiency and lightweight design. Consistent with methodologies established for DeepSeek LLMs, the optimization is carried out using the AdamW optimizer \citep{adamW}, configured with $\beta_1=0.9$, $\beta_2=0.95$, and a $\mathrm{weight\_decay}$ of $0.1$. The learning rate follows a multi-phase schedule: it linearly ramps up to its maximum over the first 2,000 warmup steps, then is reduced to 31.6\% of its maximum after reaching 80\% of total training, and ultimately lowered to 10.0\% of the initial peak by the 90\% completion mark. The maximum learning rate is specified as $5.3\text{e-}4$. A batch size comprising 4M tokens and a context window of 4K tokens are utilized for training.

\subsubsection{Evaluation Results}

\textbf{The DeepSeekMath~Corpus is of high quality, covers multilingual mathematical content, and is the largest in size.}

\begin{itemize}[topsep=0pt]
    \item \textbf{High-quality}: 
    We assess downstream capabilities across eight distinct mathematical benchmarks utilizing few-shot chain-of-thought prompting \cite{cot}. According to Table \ref{tab:corpora_comparison}, models trained with the DeepSeekMath~Corpus consistently outperform others. Additionally, as illustrated in Figure \ref{fig:corpora_comparisons}, leveraging the DeepSeekMath Corpus yields stronger results compared to Proof-Pile-2 when evaluated after processing 50B tokens (equivalent to a complete epoch over Proof-Pile-2). This underscores the superior average quality found within the DeepSeekMath Corpus.
    \item \textbf{Multilingual}:
    The DeepSeekMath~Corpus provides coverage across various languages, with English and Chinese constituting the majority of the dataset. Table \ref{tab:corpora_comparison} demonstrates that models trained on this corpus experience marked improvements in mathematical reasoning for both English and Chinese. Conversely, alternative corpora that are mainly English-focused offer minimal enhancements and, in some cases, may even impair the model's ability to reason in Chinese mathematical contexts.
    \item \textbf{Large-scale}:
    In terms of size, the DeepSeekMath~Corpus substantially exceeds existing mathematical datasets. Figure \ref{fig:corpora_comparisons} reveals that DeepSeek-LLM 1.3B, after being trained on DeepSeekMath~Corpus, achieves a more pronounced and sustained learning trajectory. On the other hand, baseline datasets, which are much smaller and repeatedly revisited during model training, lead to rapid saturation in performance.
\end{itemize}

\subsection{Training and Evaluating DeepSeekMath-Base 7B}

In this section, we present DeepSeekMath-Base 7B, a foundational model that demonstrates notable mathematical reasoning capabilities. The model’s weights are initialized with DeepSeek-Coder-Base-v1.5 7B \citep{deepseek-coder}, and it undergoes training over 500B tokens. The data composition includes 56\% sourced from the DeepSeekMath Corpus, 4\% from AlgebraicStack, 10\% from arXiv, 20\% derived from Github code, and the last 10\% made up of English and Chinese text from Common Crawl. Our training regime generally follows the configuration described in Section \ref{sec:quality-policy}, with two exceptions: the learning rate peaks at 4.2e-4, and a batch size of 10M tokens is utilized.

This work presents an in-depth evaluation of DeepSeekMath-Base 7B's proficiency in mathematics, examining its performance in generating standalone mathematical solutions independent of outside resources, tackling math problems with the aid of tools, and performing formalized theorem proving. In addition to its mathematical competencies, we further outline the broader characteristics of the base model, detailing its abilities in natural language comprehension, reasoning, and programming tasks.

\paragraph{Mathematical Problem Solving with Step-by-Step Reasoning}

We assess the problem-solving capabilities of DeepSeekMath-Base through few-shot chain-of-thought prompting \citep{cot} on eight separate benchmarks in both English and Chinese. These benchmarks span quantitative reasoning tasks—such as GSM8K \citep{gsm8k}, MATH \citep{MATH}, and CMATH \citep{wei2023cmath}—as well as multiple-choice assessments, including MMLU-STEM \citep{mmlu} and Gaokao-MathQA \citep{agieval}. Collectively, they represent a broad spectrum of mathematical disciplines, ranging in difficulty from elementary to collegiate levels.

Table \ref{tab:base_math_cot} highlights that DeepSeekMath-Base 7B consistently achieves the highest scores on all eight evaluated benchmarks when compared to other open-source base models, such as the popular general-purpose Mistral 7B \citep{mistral} and the recently introduced Llemma 34B \citep{llemma}, which was specifically trained for mathematics using Proof-Pile-2 \citep{llemma}. Particularly striking is DeepSeekMath-Base's performance on the challenging, competition-oriented MATH dataset, where it exceeds the performance of all current open-source base models by more than 10 percentage points. Furthermore, it outperforms Minerva 540B \citep{minerva}—a much larger proprietary model based on PaLM \citep{palm} and subjected to additional mathematical training—despite Minerva's scale being 77 times greater.

\paragraph{Mathematical Problem Solving with Tool Use} We assess the capability of models to perform mathematical reasoning with the assistance of tools on GSM8K and MATH benchmarks through few-shot program-of-thought prompting, following \citet{pot,pal}. For each task, models generate Python code that may leverage libraries such as \textit{math} and \textit{sympy} for conducting complex calculations. The output generated by running the code is regarded as the model’s solution. According to the results reported in Table \ref{tab:base_math_pot_proof}, DeepSeekMath-Base 7B achieves superior performance compared to the previous best, Llemma 34B.

\paragraph{Formal Mathematics}

The automation of formal proofs plays a crucial role in safeguarding the correctness and trustworthiness of mathematical reasoning, while also streamlining the proof development process. This area has received growing interest in recent years. In this study, we assess DeepSeekMath-Base 7B on the informal-to-formal proof generation task introduced by \citep{dsp_proof}, which requires producing a formal proof given an informal description, its corresponding formalized statement, and an informal argument. Our experiments make use of the miniF2F benchmark \citep{minif2f}, which targets formalizations of Olympiad-level math problems. For each item in miniF2F, we employ few-shot prompting to produce formal Isabelle proofs. In alignment with \cite{dsp_proof}, our approach involves using language models to draft proof outlines, subsequently employing the readily available Sledgehammer tool \citep{sledgehammer} to fill in gaps and complete the formal proof details. Table \ref{tab:base_math_pot_proof} presents results demonstrating that DeepSeekMath-Base 7B achieves notably effective performance in automating the formalization of proofs.

\paragraph{Natural Language Understanding, Reasoning, and Code}

We assess the natural language understanding capabilities of models on the MMLU benchmark \citep{mmlu}, reasoning proficiency using BBH \citep{bbh}, and evaluate coding skills with HumanEval \citep{codex} and MBPP \citep{mbpp}. Table \ref{tab:understanding_reasoning_code} demonstrates that DeepSeekMath-Base 7B achieves notable gains on both MMLU and BBH tasks compared to its predecessor, DeepSeek-Coder-Base-v1.5 \citep{deepseek-coder}, highlighting the beneficial effects of math-focused training on both language understanding and reasoning. Furthermore, by introducing code tokens during the phase of continual training, DeepSeekMath-Base 7B successfully preserves the coding performance levels achieved by DeepSeek-Coder-Base-v1.5 across the two relevant benchmarks. Taken together, DeepSeekMath-Base 7B substantially surpasses the general-purpose Mistral 7B model \citep{mistral} on all three examined benchmarks related to reasoning and coding.

\section{Supervised Fine-Tuning}

\subsection{SFT Data Curation}

A comprehensive mathematical instruction-tuning dataset is developed, encompassing both English and Chinese tasks drawn from a range of mathematical disciplines and spanning multiple levels of difficulty. Each problem is accompanied by solutions formatted in chain-of-thought (CoT) \citep{cot}, program-of-thought (PoT) \citep{pot,pal}, or tool-integrated reasoning structures \citep{tora}. The dataset consists of 776,000 training examples in total.

\begin{itemize}[topsep=0pt]
    \item \textbf{English mathematical datasets}: We provide annotations with tool-integrated solutions for GSM8K and MATH datasets, and incorporate a selected portion of MathInstruct \citep{MathInstruct}, as well as the Lila-OOD training data \citep{lila}, featuring solutions developed with CoT or PoT approaches. This collection of English datasets encompasses a wide range of mathematical domains, including algebra, probability, number theory, calculus, and geometry.
    \item \textbf{Chinese mathematical datasets}: Our compilation features Chinese K-12 math problems distributed across 76 sub-topics, such as linear equations, with each problem accompanied by both CoT-style and tool-integrated reasoning annotations.
\end{itemize}

\subsection{Training and Evaluating DeepSeekMath-Instruct 7B}

This section presents DeepSeekMath-Instruct 7B, developed through mathematical instruction tuning built upon DeepSeekMath-Base. During training, examples are combined in random order up to a maximum context window of 4K tokens. The model is trained across 500 iterations, employing a batch size of 256 and a fixed learning rate of 5e-5.

We evaluate models' mathematical performance both without and with tool use, on 4 quantitative reasoning benchmarks in English and Chinese. We benchmark our model against the leading models of the time:

\begin{itemize}[topsep=0pt]
    \item \textbf{Closed-source models} comprise: (1) the GPT series, where GPT-4 \citep{gpt4} and GPT-4 Code Interpreter~\footnote{\url{https://openai.com/blog/chatgpt-plugins\#\#code-interpreter}} represent the most advanced options, (2) Gemini Ultra and Pro \citep{gemini}, (3) Inflection-2 \citep{inflection-2}, (4) Grok-1~\footnote{\url{https://x.ai/model-card}}, together with several models from Chinese firms that were unveiled recently, including (5) Baichuan-3~\footnote{\url{https://www.baichuan-ai.com}}, and (6) the most recent iteration of GLM, GLM-4~\footnote{\url{https://open.bigmodel.cn/dev/api\#glm-4}}

    from the GLM family \citep{glm}.

These models are designed for broad applications and have typically been subjected to various alignment methods.  
\item \textbf{Open-source models} comprise the following: general-purpose models such as (1) DeepSeek-LLM-Chat 67B \citep{deepseek-llm}, (2) Qwen 72B \citep{qwen}, (3) SeaLLM-v2 7B \citep{seallm}, and (4) ChatGLM3 6B \citep{chatglm3}. Additionally, several models incorporate enhancements for mathematics, including (5) InternLM2-Math 20B~\footnote{\url{https://github.com/InternLM/InternLM-Math}}, which extends InternLM2 by adding mathematical training prior to instruction tuning; (6) Math-Shepherd-Mistral 7B, utilizing PPO training \citep{schulman2017proximal} on the Mistral 7B architecture \citep{mistral} through a process-supervised reward system; (7) the WizardMath collection \citep{wizardmath}, which improves the mathematical reasoning abilities of Mistral 7B and Llama-2 70B \citep{llama2} through evolve-instruct (an instruction tuning approach leveraging AI-generated prompts) and PPO training, drawing heavily on datasets such as GSM8K and MATH; (8) MetaMath 70B \citep{metamath}, where Llama-2 70B has been fine-tuned on an expanded version of GSM8K and MATH; (9) ToRA 34B \cite{tora}, which adapts CodeLlama 34B for mathematical reasoning that integrates external tools; and (10) MAmmoTH 70B \citep{MathInstruct}, a Llama-2 70B variant that is instruction-tuned with the MathInstruct dataset.  
\end{itemize}

Referencing the results in Table \ref{tab:sft_rl_math}, when tool assistance is restricted, DeepSeekMath-Instruct 7B achieves robust step-by-step reasoning performance. In particular, on the competition-grade MATH dataset, this model outperforms all other open-source alternatives as well as most proprietary systems (including Inflection-2 and Gemini Pro) by at least a 9\% margin in absolute terms. This advantage persists even in comparison to significantly larger models (such as Qwen 72B) or those benefitting from dedicated math-oriented reinforcement learning strategies (for instance, WizardMath-v1.1 7B). Although DeepSeekMath-Instruct's performance approaches that of Chinese proprietary models like GLM-4 and Baichuan-3 on the MATH benchmark, it continues to lag behind the capabilities of GPT-4 and Gemini Ultra.

When evaluated under conditions that permit models to combine both natural language reasoning and programmatic tool utilization for solving problems, DeepSeekMath-Instruct 7B achieves close to 60\% accuracy on MATH, outperforming all previously released open-source models. Across additional benchmark datasets, our model demonstrates performance on par with DeepSeek-LLM-Chat 67B, which previously held the state-of-the-art status despite being ten times the size.

\subsection{Group Relative Policy Optimization} After the Supervised Fine-Tuning (SFT) phase, reinforcement learning (RL) has demonstrated its utility in enhancing the mathematical reasoning skills of LLMs \citep{wang2023math,wizardmath}. In the following, we present our novel RL technique, termed Group Relative Policy Optimization (GRPO), which is both efficient and effective.

\subsubsection{From PPO to GRPO} Proximal Policy Optimization (PPO) \citep{schulman2017proximal} is an actor-critic RL algorithm that is widely used in the RL fine-tuning stage of LLMs \citep{ouyang2022training}. In particular, it optimizes LLMs by maximizing the following surrogate objective:

\begin{equation}
\footnotesize
    \mathcal{J}_{PPO}(\theta) = \mathbb{E}{[q \sim P(Q), o \sim \pi_{\theta_{old}}(O|q)]} \frac{1}{|o|} \sum_{t=1}^{|o|} \min \left[ \frac{\pi_\theta(o_{t} | q, o_{<t})}{\pi_{\theta_{old}}(o_{t} | q, o_{<t})} A_{t}, \text{clip} \left( \frac{\pi_\theta(o_{t} | q, o_{<t})}{\pi_{\theta_{old}}(o_{t} | q, o_{<t})}, 1 - \varepsilon, 1 + \varepsilon \right)  A_{t} \right] ,
\end{equation}

The equation above represents the PPO objective, where expectations are computed over query-sample pairs drawn according to $q \sim P(Q)$ and $o \sim \pi_{\theta_{old}}(O|q)$. The loss aggregates over all positions $t$ within the output sequence, normalizing by the sequence length $|o|$. For each position, the minimum is calculated between the product of the advantage function $A_t$ with the ratio of current to previous policy probabilities for $o_t$, and the clipped surrogate term where the ratio is limited within $1-\varepsilon$ and $1+\varepsilon$ before being multiplied by $A_t$.

Here, $\pi_{\theta}$ denotes the present policy model, while $\pi_{\theta_{old}}$ refers to the preceding one; $q$ and $o$ represent questions drawn from the question dataset and outputs generated by the previous policy $\pi_{\theta_{old}}$, respectively. The parameter $\varepsilon$ serves as a PPO-specific clipping threshold, playing a key role in maintaining training stability. The term $A_t$ stands for the advantage, which is estimated using Generalized Advantage Estimation (GAE) \citep{gae} based on the reward sequence $\{r_{\ge t}\}$ and a trainable value function $V_{\psi}$. Consequently, PPO necessitates concurrent training of a value function alongside the main policy. Additionally, to prevent the reward model from being excessively optimized, a typical strategy involves introducing a KL divergence penalty—applied per token and measured against a reference model—within the reward structure at each token, as detailed in \citep{ouyang2022training}, that is,

\begin{equation}
    r_{t} = r_\varphi(q, o_{\le t}) - \beta \log\frac{\pi_{\theta}(o_{t}|q, o_{<t})}{\pi_{ref}(o_{t}|q, o_{<t})},
\label{eq:PPO-reward}
\end{equation}

Here, $r_\varphi$ denotes the reward model, $\pi_{ref}$ refers to the reference model—commonly the original SFT model—and $\beta$ represents the KL penalty coefficient.

As the value function employed in PPO is typically another model of comparable size as the policy model, it brings a substantial memory and computational burden. Additionally, during RL training, the value function is treated as a baseline in the calculation of the advantage for variance reduction. While in the LLM context, usually only the last token is assigned a reward score by the reward model, which may complicate the training of a value function that is accurate at each token. To address this, as shown in Figure \ref{fig:grpo}, we propose Group Relative Policy Optimization (GRPO), which obviates the need for additional value function approximation as in PPO, and instead uses the average reward of multiple sampled outputs, produced in response to the same question, as the baseline. More specifically, for each question $q$, GRPO samples a group of outputs $\{o_1, o_2, \cdots, o_G\}$  from the old policy  $\pi_{\theta_{old}}$  and then optimizes the policy model by maximizing the following objective:

\begin{equation}
\footnotesize
\begin{split}
    \mathcal{J}_{GRPO}(\theta) &= \mathbb{E}{[q \sim P(Q), \{o_i\}_{i=1}^G \sim \pi_{\theta_{old}}(O|q)]} \\
    & \frac{1}{G}\sum_{i=1}^G\frac{1}{|o_i|} \sum_{t=1}^{|o_i|} \left\{ \min \left[ \frac{\pi_\theta(o_{i,t} | q, o_{i,<t})}{\pi_{\theta_{old}}(o_{i,t} | q, o_{i,<t})} \hat{A}_{i,t}, \text{clip} \left( \frac{\pi_\theta(o_{i,t} | q, o_{i,<t})}{\pi_{\theta_{old}}(o_{i,t} | q, o_{i,<t})}, 1 - \varepsilon, 1 + \varepsilon \right)  \hat{A}_{i,t} \right] - \beta \mathbb{D}_{KL}\left[\pi_{\theta} || \pi_{ref}\right]\right\} ,
\end{split}
\label{eq:GRPO-obj}
\end{equation}

where $\varepsilon$ and $\beta$ are hyper-parameters, and $\hat{A}_{i,t}$  is the advantage calculated based on relative rewards of the outputs inside each group only, which will be detailed in the following subsections. The group relative way that GRPO leverages to calculate the advantages, aligns well with the comparative nature of rewards models,  as reward models are typically trained on datasets of comparisons between outputs on the same question. Also note that, instead of adding KL penalty in the reward, GRPO regularizes by directly adding the KL divergence between the trained policy and the reference policy to the loss, avoiding complicating the calculation of  $\hat{A}_{i,t}$. And different from the KL penalty term used in (\ref{eq:PPO-reward}), we estimate the KL divergence with the following unbiased estimator \citep{kl_approx}:

\begin{equation}
\small
    \mathbb{D}_{KL}\left[\pi_{\theta} || \pi_{ref}\right] = \frac{\pi_{ref}(o_{i,t}|q,o_{i,<t})}{\pi_{\theta}(o_{i,t}|q,o_{i,<t})} - \log\frac{\pi_{ref}(o_{i,t}|q,o_{i,<t})}{\pi_{\theta}(o_{i,t}|q,o_{i,<t})} - 1,
\end{equation}

which is guaranteed to be positive.

\begin{algorithm}[t]
  \small
  \caption{Iterative Group Relative Policy Optimization}
  \textbf{Input} initial policy model $\pi_{\theta_{\text{init}}}$; reward models $r_\varphi$; task prompts $\mathcal{D}$; 
  hyperparameters $\varepsilon$, $\beta$, $\mu$
  \begin{algorithmic}[1]
    \State Set policy model $\pi_\theta$ to $\pi_{\theta_{\text{init}}}$
    \For{iteration = 1 \textbf{to} I}
      \State Assign reference model $\pi_{ref} \gets \pi_\theta$
      \For{step = 1 \textbf{to} M}
        \State Select a batch $\mathcal{D}_b$ from $\mathcal{D}$
        \State Store previous policy model $\pi_{\theta_{old}} \gets \pi_\theta$
        \State Generate $G$ outputs $\{o_i\}_{i=1}^G$ according to $\pi_{\theta_{old}} (\cdot \mid q)$ for each question $q$ in $\mathcal{D}_b$
        \State Evaluate and obtain rewards $\{r_i\}_{i=1}^{G}$ for each $o_i$ by applying $r_\varphi$
        \State Estimate group relative advantage $\hat{A}_{i,t}$ for token $t$ within each $o_i$
        \For{GRPO iteration = 1 \textbf{to} $\mu$}
            \State Update $\pi_{\theta}$ by maximizing the GRPO objective specified by Equation \ref{eq:GC-GRPO}
        \EndFor
      \EndFor
      \State Continually optimize $r_\varphi$ with a replay buffer.
    \EndFor
  \end{algorithmic}
  \textbf{Output} $\pi_\theta$
  \label{alg:iter-grpo}
\end{algorithm}

\subsubsection{Outcome Supervision RL with GRPO} In this approach, given a question $q$, a collection of responses $\{o_1, o_2, \cdots, o_G\}$ is drawn using the prior policy model $\pi_{\theta_{old}}$. Each response is then assigned a reward by a dedicated reward model, resulting in a reward set $\mathbf{r}=\{r_1, r_2, \cdots, r_G\}$. These reward values are then standardized within the group by removing the mean and scaling by the standard deviation of $\mathbf{r}$. Outcome supervision delivers this normalized reward to the terminal step of each output $o_i$, while also using the normalized reward to define the advantage $\hat{A}_{i, t}$ for every token position within $o_i$, expressed as $\hat{A}_{i, t} = \widetilde{r}_i = \frac{r_i- {\rm mean}(\mathbf{r})}{{\rm std}(\mathbf{r})}$. The policy is then updated through maximization of the objective specified in equation (\ref{eq:GRPO-obj}).

\subsubsection{Process Supervision RL with GRPO} Traditional outcome supervision delivers feedback solely at the conclusion of each generated output, which often proves inadequate and inefficient when guiding policies on complex mathematical problems. Inspired by \cite{wang2023math}, our work incorporates process supervision, wherein feedback is assigned after each intermediate reasoning step. Specifically, for a given question $q$ and a collection of $G$ sampled outputs $\{o_1, o_2, \cdots, o_G\}$, a dedicated process reward model evaluates every step within the outputs, producing reward sets: $\mathbf{R} = \{ \{r_1^{index(1)},\cdots,r_1^{index(K_1)}\}, \cdots,  \{r_G^{index(1)},\cdots,r_G^{index(K_G)}\} \}$. Here, $index(j)$ denotes the end token of the $j$-th reasoning step, and $K_i$ captures the number of steps within output $i$. These step-level rewards undergo normalization using their mean and standard deviation, yielding: $\widetilde{r}_i^{index(j)} = \frac{r_i^{index(j)} - {\rm mean(\mathbf{R})}}{{\rm std(\mathbf{R})}}$.
For each token, we then compute its advantage by summing the normalized rewards across all subsequent steps, formally given by $\hat{A}_{i, t} = \sum_{index(j) \ge t} \widetilde{r}_i^{index(j)}$. This advantage is employed to update the policy by maximizing the objective shown in equation (\ref{eq:GRPO-obj}).

\subsubsection{Iterative RL with GRPO} During reinforcement learning, the previously trained reward model may become inadequate for effectively guiding the increasingly sophisticated policy model. To address this, we investigate iterative RL with GRPO. As detailed in Algorithm \ref{alg:iter-grpo}, the iterative GRPO procedure entails creating fresh datasets for reward model training using new samples obtained from the policy model. The reward model is further updated using a replay buffer containing 10\% of earlier data. Subsequently, the policy model serves as the new reference model, and training continues using the updated reward model to further refine policy performance.

\subsection{Training and Evaluating DeepSeekMath-RL}

We implement reinforcement learning using DeepSeekMath-Instruct 7B as our base. The RL training leverages chain-of-thought formatted questions, specifically selected from GSM8K and MATH within the SFT dataset, totaling approximately 144K examples. To specifically examine RL's effects on benchmarks that remain unseen during RL, all other SFT questions are excluded from this process. The reward model training set is assembled in line with the procedure described in \citep{wang2023math}. Our initial reward model employs DeepSeekMath-Base 7B, trained with a learning rate of $2 \times 10^{-5}$. In the case of GRPO, the policy model utilizes a learning rate of $1 \times 10^{-6}$. A KL penalty coefficient of $0.04$ is applied. For each prompt, we generate $64$ sampled outputs. The maximum sequence length permitted during training is 1024 tokens, with a batch size of 1024. After each exploration period, the policy model receives only one update. DeepSeekMath-RL 7B is evaluated using the same suite of benchmarks as DeepSeekMath-Instruct 7B. Specifically, for DeepSeekMath-RL 7B, tasks involving GSM8K and MATH questions with chain-of-thought reasoning are treated as in-domain, whereas all remaining benchmarks are considered out-of-domain.

Table \ref{tab:sft_rl_math} presents the results for both open-source and closed-source models using chain-of-thought and tool-augmented reasoning on English and Chinese evaluation tasks. The findings can be summarized as follows: (1) Employing chain-of-thought reasoning, DeepSeekMath-RL 7B achieves accuracy rates of 88.2\% on GSM8K and 51.7\% on MATH. These results not only exceed those of all open-source models within the 7B–70B parameter range, but also outperform most closed-source counterparts. (2) Notably, the training of DeepSeekMath-RL 7B relies exclusively on chain-of-thought-based instruction tuning datasets from GSM8K and MATH, with DeepSeekMath-Instruct 7B as the initialization point. Even with this limited training corpus, the model demonstrates consistent improvements over DeepSeekMath-Instruct 7B in every evaluation metric, thereby illustrating the strengths of reinforcement learning in this context.

\section{Discussion}
This section presents our observations from both pre-training and reinforcement learning experiments.

\subsection{Lessons Learnt in Pre-Training}

We begin by discussing our observations during the pre-training phase. Unless explicitly mentioned, our methodology follows the training configurations described in Section \ref{sec:quality-policy}. Additionally, it should be emphasized that references to the DeepSeekMath Corpus in this section pertain to an 89B-token subset obtained from the second round of data collection.

\subsubsection{Code Training Benefits Mathematical Reasoning}

An often-cited but largely untested hypothesis is that training on code can enhance reasoning capabilities. In this work, we seek to partially address this claim, focusing specifically on mathematical reasoning: code training enhances models' proficiency in mathematical reasoning tasks, regardless of whether external tools are employed.

To study how code training affects mathematical reasoning, we experimented with the following two-stage training and one-stage training settings:

\textbf{Two-Stage Training}

\begin{itemize}[topsep=0pt]
    \item \textbf{Code Training for 400B Tokens $\rightarrow$ Math Training for 150B Tokens}: Our approach involves first exposing DeepSeek-LLM 1.3B to 400B code tokens, subsequently followed by an additional 150B tokens focused on mathematical content;
    \item \textbf{General Training for 400B Tokens $\rightarrow$ Math Training for 150B Tokens}: For comparison, we perform a baseline experiment where, in the initial phase, the model is trained with 400B tokens selected from a broad general-purpose corpus developed by DeepSeek-AI instead of code tokens, aiming to assess whether code token pre-training provides unique benefits for enhancing mathematical reasoning capabilities.
\end{itemize}

\textbf{One-Stage Training}

\begin{itemize}[topsep=0pt]
    \item \textbf{150B Token Math Pretraining}: DeepSeek-LLM 1.3B undergoes pretraining on a dataset containing 150B math tokens;
    \item \textbf{Combined Training with 400B Code Tokens and 150B Math Tokens}: Sequentially training on math after code has been found to reduce coding abilities. We therefore explore if simultaneous, one-stage training on both math and code token streams not only enhances mathematical reasoning but also helps counteract the issue of catastrophic forgetting in code proficiency.
\end{itemize}

\paragraph{Results}
Tables \ref{tab:code-math} and \ref{tab:code-math-general-eval} present the downstream outcomes across various training configurations.

Code training enhances program-aided mathematical reasoning across both two-stage and one-stage training paradigms. According to Table \ref{tab:code-math}, in the two-stage setup, exclusive code training substantially improves performance on GSM8K and MATH tasks utilizing Python. Subsequent math-specific training during the second stage leads to additional gains. Notably, in the one-stage setting, jointly training with both code and math tokens effectively addresses the problem of catastrophic forgetting commonly seen with two-stage approaches, while also delivering improvements in coding abilities (Table \ref{tab:code-math-general-eval}) and program-assisted mathematical reasoning (Table \ref{tab:code-math}).

Engaging in code training likewise strengthens mathematical reasoning abilities in the absence of external tools. In the context of two-stage training, undertaking code training as the first phase yields noticeable improvements early on. This process also facilitates more effective math training in the later stage, culminating in optimal results. On the other hand, merging code tokens and math tokens during a single-phase training regimen diminishes mathematical reasoning performance when tools are not employed. A possible explanation is that, given the relatively small size of DeepSeek-LLM 1.3B, the model may not possess adequate capacity to thoroughly learn and integrate both code and math information within a unified training stage.

\subsubsection{ArXiv Papers Seem Ineffective in Improving Mathematical Reasoning}

ArXiv papers are commonly included as a component of math pre-training data \citep{minerva,gpt-f,llemma,mathpile}. However, detailed analysis regarding their impact on mathematical reasoning has not been extensively conducted. Perhaps counter-intuitively, according to our experiments, arXiv papers seem ineffective in improving mathematical reasoning. We experiment with models of different sizes, including DeepSeek-LLM 1.3B and DeepSeek-Coder-Base-v1.5 7B \citep{deepseek-coder}, using arXiv corpora that underwent varied processing pipelines:

\begin{itemize}[topsep=0pt]
    \item \textbf{MathPile} \citep{mathpile}: a dataset comprising 8.9 billion tokens, constructed through a series of data cleaning and filtering strategies, with scientific arXiv documents constituting over 85\% of its content;
    \item \textbf{ArXiv-RedPajama} \citep{redpajama}: a collection containing all arXiv LaTeX source files, from which elements such as preambles, comments, macros, and bibliographies have been excised, amounting to 28.0 billion tokens in total.
\end{itemize}

Within the scope of our experiments, DeepSeek-LLM 1.3B was individually trained on 150B tokens and DeepSeek-Coder-Base-v1.5 7B was trained on 40B tokens, each using an arXiv-specific corpus. Our findings indicate that arXiv papers provide little to no benefit in advancing models' mathematical reasoning abilities. When these models are exposed exclusively to arXiv data, performance remains stagnant or even declines on a range of mathematical benchmarks with varying levels of complexity examined in this work. These benchmarks encompass quantitative reasoning datasets such as GSM8K and MATH (Table \ref{tab:arxiv-cot}), multiple-choice assessments like MMLU-STEM (Table \ref{tab:arxiv-cot}), as well as formal mathematics tasks, including miniF2F (Table \ref{tab:arxiv_atp}).

However, this conclusion has its limitations and should be taken with a grain of salt. We have not yet studied:

\begin{itemize}[topsep=0pt]
    \item How arXiv tokens influence mathematical tasks beyond those examined here, including tasks like the informalization of theorems, which involves transforming formal mathematical statements or proofs into their informal counterparts;
    \item The consequences of integrating arXiv tokens with other varieties of data;
    \item If the advantages observed from utilizing arXiv papers would persist or become more evident with scaling to larger models.
\end{itemize}

Thus, further exploration is required, which we leave for future studies.

\subsection{Insights of Reinforcement Learning}
\subsubsection{Towards to a Unified Paradigm} Here, we introduce a cohesive framework for evaluating various training approaches, including SFT, RFT, DPO, PPO, and GRPO. Furthermore, we undertake empirical studies to investigate key components influencing this unified paradigm. In a general sense, the gradient with respect to the parameter $\theta$ for any given training algorithm can be formulated as follows:

\begin{equation}
    \nabla_{\theta}\mathcal{J}_{\textcolor{red}{\mathcal{A}}}(\theta) = \mathbb{E}[\underbrace{(q,o) \sim \textcolor{red}{\mathcal{D}}}_{Data \ Source}]\left( \frac{1}{|o|} \sum_{t=1}^{|o|}  \underbrace{GC_{{\mathcal{A}}}(q, o, t, \textcolor{red}{\pi_{{rf}}})}_{Gradient \ Coefficient}  \nabla_{\theta}\log \pi_{\theta}(o_t | q, o_{<t})\right).
\label{eq:objective}
\end{equation}

There exist three key components: 1) \textit{Data Source $\mathcal{D}$}, which determines the training data; 2) \textit{Reward Function $\pi_{{rf}}$}, which is the source of the training reward signal; 3) \textit{Algorithm $\mathcal{A}$}: which processes the training data and the reward signal to the gradient coefficient $GC$ that determines the magnitude of the penalty or reinforcement for the data. We analyze several representative methods based on such a unified paradigm:

\begin{itemize}
    \item  \textbf{Supervised Fine-tuning (SFT)}: SFT adapts a pretrained model by further training it on datasets curated and labeled by humans for fine-tuning purposes.
    \item \textbf{Rejection Sampling Fine-tuning (RFT)}: RFT builds upon the SFT model by additional training on a subset of model outputs generated in response to SFT prompts, selectively retaining only those outputs assessed as correct for further learning.
    \item \textbf{Direct Preference Optimization (DPO)}: DPO enhances the SFT model by fine-tuning it on a set of model-generated responses, leveraging a pair-wise DPO loss to incorporate preferences between alternative outputs.
    \item \textbf{Online Rejection Sampling Fine-tuning (Online RFT)}: In contrast to standard RFT, Online RFT begins with the SFT model as the initial policy and updates it by fine-tuning using newly generated outputs sampled in real time from the evolving policy model.
    \item \textbf{PPO/GRPO}: PPO/GRPO starts from the SFT-trained policy model and employs reinforcement learning, continuously improving it with outputs generated on-the-fly by the current version of the policy.
\end{itemize}

Table \ref{tab:data_policy} provides an overview of the elements involved in these approaches. For a comprehensive explanation of the derivation process, see Appendix \ref{app:analysis-rl}.

\paragraph{Observation about Data Source} We classify data sources into two main types: online sampling and offline sampling. In online sampling, training data is derived from the exploratory behavior of the current, actively trained policy model. In contrast, offline sampling refers to training data collected from outputs generated by the original SFT model. RFT and DPO utilize the offline sampling approach, whereas Online RFT and GRPO adopt the online sampling methodology.

As illustrated in Figure \ref{fig:rl-analysis}, our results indicate that Online RFT achieves markedly better performance than RFT across both benchmarks. While the performance of Online RFT closely matches that of RFT during the initial training phase, it consistently surpasses RFT in the latter stages, underscoring the benefits of online learning. This outcome aligns with expectations: early in training, the behavior of the actor is still very similar to that of the SFT model, resulting in sampled data with only subtle distinctions. As training progresses, the differences in samples drawn from the actor become more pronounced, making real-time sampling increasingly advantageous.

\paragraph{Observation about Gradient Coefficient} In our approach, the transformation of the reward signal into a gradient coefficient is central to updating model parameters. For the purposes of our study, we distinguish the reward function into two types: ‘Rule’ and ‘Model’. ‘Rule’ refers to assessments grounded in the objective correctness of an answer, whereas ‘Model’ indicates the use of a trained reward model that scores responses; the latter’s training data is sourced from rule-based judgments. Examining Equations \ref{eq:GC-RFT} and \ref{eq:GC-GRPO}, it becomes evident that GRPO differs from Online RFT in a fundamental way: the gradient coefficient in GRPO dynamically changes in accordance with the magnitude of rewards generated by the reward model. This mechanism enables variable degrees of positive reinforcement or penalization, based on individual response scores. By contrast, Online RFT does not employ this flexibility; it applies no penalties to incorrect outputs and always provides a uniform level of positive reinforcement to all correct responses, regardless of their quality.

As illustrated in Figure \ref{fig:rl-analysis}, GRPO achieves better results than online RFT, underscoring the advantage of modifying the positive and negative gradient coefficients. Moreover, the combination of GRPO+PS consistently outperforms GRPO+OS, which suggests that employing detailed, step-sensitive gradient coefficients is beneficial. Additionally, we investigate the effect of iterative RL by performing two successive iterations in our experimental setup. Figure \ref{fig:iter-rl} reveals that iterative RL yields marked performance gains, with the most substantial improvement observed during the initial iteration.

\subsubsection{Why RL Works?} In this study, we implement reinforcement learning using a selected subset of instruction tuning data, resulting in considerable improvements over the baseline instruction tuning model. To better understand the effectiveness of reinforcement learning, we measure both Pass@K and Maj@K accuracies for the Instruct and RL models on two different benchmarks. The results presented in Figure \ref{fig:maj-pass} indicate that RL leads to improvements in Maj@K but does not affect Pass@K. This pattern suggests that reinforcement learning strengthens the robustness of the output distribution, namely, \textbf{the observed gains may primarily stem from increasing the chance of obtaining the correct answer among the TopK outputs instead of fundamentally advancing the model’s intrinsic skills.} Likewise, \citep{wang2023making} observed a \textbf{misalignment problem} in SFT models applied to reasoning tasks, and demonstrated that the reasoning capabilities of SFT models could be enhanced by employing preference alignment strategies \citep{yuan2023rrhf,song2023preference,wang2023making}.

\subsubsection{How to Achieve More Effective RL?}
\label{sec:effec_rl}
Our results show that RL is quite effective when applied to mathematical reasoning challenges. Furthermore, we introduce a unified framework that facilitates the interpretation of various prominent training strategies, viewing them all as forms of direct or reduced RL methods. As highlighted in Equation \ref{eq:objective}, the core elements involved include the Data Source, Algorithm, and Reward Function. We suggest several possible avenues for future investigation pertaining to these three fundamental aspects.

\paragraph{Data Source} The data source serves as the foundational input for all forms of training. Within the realm of RL, we define the data source as unlabeled questions paired with outputs generated by sampling from the policy model. For the experiments in this work, our approach relies exclusively on questions obtained during the instruction tuning phase, and employs a straightforward nucleus sampling strategy for generating outputs. We hypothesize that this limitation may partly account for our RL pipeline’s improvements being confined to the Maj@K metric. Future work will extend our RL pipeline to out-of-distribution prompt scenarios, potentially integrating \textbf{advanced sampling (decoding) strategies}, such as those utilizing tree-search methods \citep{yao2023tree}. Additionally, the development of \textbf{efficient inference techniques} \citep{xia-etal-2023-speculative,leviathan2023fast,kwon2023efficient,xia2024unlocking}, which directly impact the exploratory capabilities of policy models, remains a critical avenue of investigation.

\paragraph{Algorithms} Algorithms utilize both data and reward signals to compute the gradient coefficient, which in turn updates model parameters. In accordance with Equation \ref{eq:objective}, contemporary approaches fundamentally \textbf{TRUST} the reward function's signal to either elevate or suppress the conditional probability associated with particular tokens. Nonetheless, the assumption that the reward signal is consistently dependable does not hold, particularly when confronting highly intricate tasks. For instance, the PRM800K datasets \citep{lightman2023let}, although meticulously labeled by expert annotators, still possess an error rate of about 20\% in their annotations\footnote{\url{https://github.com/openai/prm800k/issues/12\#issuecomment-1728491852}}. In light of this, our focus shifts towards investigating reinforcement learning algorithms capable of withstanding noisy reward feedback. We anticipate that \textbf{WEAK-TO-STRONG} alignment strategies \citep{burns2023weak} can introduce transformative advancements to the underlying learning frameworks.

\paragraph{Reward Function} The reward function serves as the principal source of learning signals during training. In the context of RL, it is typically instantiated as a neural reward model. We identify three pivotal avenues for advancing reward models: 1) \textbf{Strategies for bolstering the generalization capacity of the reward model.} For reinforcement learning to genuinely enhance the intrinsic abilities of large language models (LLMs), the reward model needs to reliably generalize to novel, out-of-distribution queries and sophisticated decoding sequences; absent this, RL may merely help maintain the existing output distribution rather than drive substantive improvements; 2) \textbf{Incorporating representations of the reward model's uncertainty.} Capturing uncertainty may serve as a critical interface between underperforming reward models and algorithms that facilitate weak-to-strong learning; 3) \textbf{Development of efficient, high-quality process reward models} capable of supplying nuanced, fine-grained feedback during the stepwise reasoning process \citep{lightman2023let,wang2023math}.

\section{Conclusion, Limitation, and Future Work}

In this work, we introduce DeepSeekMath, a model that achieves superior results compared to all other open-source systems on the competition-level MATH benchmark and nearly matches the scores of proprietary models. DeepSeekMath~builds upon DeepSeek-Coder-v1.5 7B, undergoing extended continual pretraining over 500B tokens. Notably, 120B of these tokens are math-focused, curated from Common Crawl data. Through a comprehensive ablation analysis, we determine that web pages represent a rich reservoir for high-quality mathematical content, whereas arXiv data does not provide the expected benefits. Furthermore, we propose Group Relative Policy Optimization (GRPO), a new variant of Proximal Policy Optimization (PPO), which enhances mathematical reasoning performance while using less memory. Experimental evidence suggests that GRPO maintains its effectiveness, even as DeepSeekMath-Instruct 7B achieves elevated benchmark results. Lastly, we offer a unified framework for interpreting various approaches and highlight promising avenues for advancing reinforcement learning methods in this context.

While DeepSeekMath~demonstrates strong results on quantitative reasoning tasks, its performance in geometry-related and theorem-proving tasks is comparatively lower than that of proprietary models. For example, in our preliminary evaluation, the model was unable to solve certain problems involving triangles and ellipses, suggesting a possible bias in data selection during pre-training and fine-tuning phases. Furthermore, DeepSeekMath~exhibits inferior few-shot learning abilities relative to GPT-4, likely due to its model size limitations. Unlike GPT-4, which benefits from few-shot inputs to boost its accuracy, DeepSeekMath~displays little difference between its zero-shot and few-shot results. Moving forward, we aim to enhance our curated data selection process to assemble a higher-quality pre-training dataset. Additionally, we plan to investigate novel approaches for reinforcement learning of LLMs, as discussed in Section~\ref{sec:effec_rl}.

\end{document}